\documentclass[a4paper]{article}
\usepackage[a4paper,margin=0pt]{geometry}
\usepackage{ctex}
\usepackage{amsmath,amssymb}
\usepackage{array,booktabs,graphicx}
\usepackage{xcolor}
\usepackage{tikz}
\usetikzlibrary{arrows.meta,calc,positioning}
\pagestyle{empty}
\setlength{\parindent}{0pt}
\setlength{\tabcolsep}{1.4pt}
\renewcommand{\arraystretch}{0.88}
\graphicspath{{./assets/}{course-content/runtime/lessons/review/assets/}}

\definecolor{sheetblue}{RGB}{0,82,168}
\definecolor{ruleblue}{RGB}{40,102,180}
\newcommand{\boxfont}{\fontsize{5.35pt}{5.95pt}\selectfont}
\newcommand{\headfont}{\fontsize{8pt}{8.6pt}\selectfont\bfseries}
\newcommand{\tinyformula}{\fontsize{5pt}{5.7pt}\selectfont}
\setlength{\fboxsep}{.9mm}
\setlength{\fboxrule}{.22pt}
\newcommand{\drawbox}[5]{%
  \draw[ruleblue,line width=.35pt] (#1,#2) rectangle ++(#3,#4);
  \fill[sheetblue] (#1,#2) rectangle ++(#3,4.5);
  \node[white,font=\headfont,anchor=center] at (#1+#3/2,#2+2.25) {#5};
}
\newcommand{\contentnode}[4]{%
  \node[anchor=north west,inner sep=.75mm,text width=#3,align=left] at (#1,#2) {\boxfont #4};
}
\newcommand{\subh}[1]{\noindent{\color{sheetblue}\bfseries #1}\par\vspace{.15mm}}
\newcommand{\subbox}[2]{%
  \noindent\fcolorbox{ruleblue}{white}{%
    \begin{minipage}[t]{#1}
      \boxfont #2
    \end{minipage}}%
  \par\vspace{.35mm}%
}

\begin{document}
\begin{tikzpicture}[x=1mm,y=-1mm]
  \draw[black,line width=.45pt] (0.5,0.5) rectangle (209.5,296.5);

  \node[anchor=north,font=\fontsize{15pt}{16pt}\selectfont\bfseries] at (105,1.2)
    {经典控制理论 CHEAT SHEET};
  \node[anchor=north,font=\fontsize{7.5pt}{8pt}\selectfont\bfseries,text=sheetblue] at (105,10.7)
    {建模 \quad | \quad 时域分析 \quad | \quad 根轨迹 \quad | \quad 频域分析 \quad | \quad 典型校正方法};

  \draw[ruleblue,line width=.35pt] (1.5,17.2) rectangle (208.5,28.5);
  \node[anchor=north west,font=\fontsize{5.2pt}{5.8pt}\selectfont\bfseries] at (3,18.1)
    {常用符号：};
  \node[anchor=north west,font=\fontsize{4.95pt}{5.45pt}\selectfont,text width=188mm,align=left,inner sep=0] at (18,18.1)
    {$r(t)$ 输入 \quad $c(t)$ 输出 \quad $e(t)=r(t)-b(t)$ 误差 \quad
     $G(s)$ 前向通道 \quad $H(s)$ 反馈通道 \quad $L(s)=G(s)H(s)$ 开环传递函数 \quad
     $T(s)=C(s)/R(s)$ 闭环传递函数\\[-.3mm]
     $p$ 极点 \quad $z$ 零点 \quad $\omega$ 角频率(rad/s) \quad $\zeta$ 阻尼比
     \quad $\omega_n$ 自然频率 \quad $\omega_d$ 阻尼振荡频率 \quad $K$ 增益 \quad $\tau$ 时间常数/延迟};

  \drawbox{1.5}{30}{67.5}{146}{1. 建模基础}
  \drawbox{70.5}{30}{68}{146}{2. 时域分析}
  \drawbox{140}{30}{68.5}{88}{3. 根轨迹分析}
  \drawbox{140}{119}{68.5}{77}{4. 频域分析}
  \drawbox{1.5}{177}{67.5}{87}{5. 典型校正方法}
  \drawbox{70.5}{177}{68}{87}{6. 经典结论速查}
  \drawbox{140}{197}{68.5}{67}{4.6-4.8 稳定裕度}
  \drawbox{1.5}{265}{207}{26}{7. 图形速览}

  \contentnode{2.4}{35.2}{65mm}{
    \subbox{61.6mm}{\subh{1.1 标准负反馈结构}
    \begin{tikzpicture}[x=1mm,y=1mm,>=Stealth,line width=.45pt,baseline=(current bounding box.north)]
      \node (r) at (0,10) {$R(s)$};
      \node[circle,draw,minimum size=5mm,inner sep=0] (sum) at (11,10) {};
      \node[font=\tiny] at (9.3,12.4) {$+$};
      \node[font=\tiny] at (11,6.9) {$-$};
      \node[draw,minimum width=10mm,minimum height=5mm] (g) at (25,10) {$G(s)$};
      \node (c) at (43,10) {$C(s)$};
      \node[draw,minimum width=10mm,minimum height=5mm] (h) at (25,0) {$H(s)$};
      \draw[->] (r)--(sum); \draw[->] (sum)--(g); \draw[->] (g)--(c);
      \draw[->] (39,10)|-(h); \draw[->] (h)-|(sum);
    \end{tikzpicture}
    \[
      T(s)=\frac{C(s)}{R(s)}=\frac{G(s)}{1+G(s)H(s)},\quad
      \frac{E(s)}{R(s)}=\frac{1}{1+G(s)H(s)}
    \]}
    \subbox{61.6mm}{\subh{微分方程与传递函数}
    \[
      \sum a_i y^{(i)}=\sum b_j u^{(j)},\qquad
      G(s)=\frac{Y(s)}{U(s)}
    \]}
    \subbox{61.6mm}{\subh{1.2 基本环节模型}
    \begin{tabular}{|c|c|c|c|c|c|}
      \hline
      比例 & 惯性 & 积分 & 微分 & 二阶 & 滞后\\ \hline
      $K$ & $\frac{K}{Ts+1}$ & $\frac{K}{s}$ & $Ks$ &
      $\frac{\omega_n^2}{s^2+2\zeta\omega_ns+\omega_n^2}$ & $e^{-\tau s}$\\ \hline
    \end{tabular}
    }
    \subbox{61.6mm}{\subh{1.3 物理系统建模}
    \begin{tabular}{|c|c|c|}
      \hline
      元件 & 关系式 & 复频域\\ \hline
      质量$M$ & $f=M\ddot x$ & $X/F=1/(Ms^2)$\\
      阻尼$B$ & $f=B\dot x$ & $X/F=1/(Bs)$\\
      弹簧$K$ & $f=Kx$ & $X/F=1/K$\\
      电阻$R$ & $v=Ri$ & $R$\\
      电感$L$ & $v=L\dot i$ & $Ls$\\
      电容$C$ & $i=C\dot v$ & $1/(Cs)$\\ \hline
    \end{tabular}
    }
    \subbox{61.6mm}{\subh{1.4 方框图代数与梅森公式}
    \[
      G=G_1G_2,\quad G=G_1+G_2,\quad T=\frac{G}{1\pm GH}
    \]
    \[
      G=\frac{\sum_{k=1}^N P_k\Delta_k}{\Delta},\qquad
      \Delta=1-\sum L_i+\sum L_iL_j-\cdots
    \]
    $P_k$ 为第 $k$ 条前向通路增益；$\Delta_k$ 为去除与该通路接触回路后的特征式。
    }
  }

  \contentnode{71.4}{35.2}{66mm}{
    \subbox{62.1mm}{\subh{2.1 标准二阶系统}
    \[
      \Phi(s)=\frac{\omega_n^2}{s^2+2\zeta\omega_ns+\omega_n^2},\quad
      \omega_d=\omega_n\sqrt{1-\zeta^2},\quad
      \sigma=\zeta\omega_n
    \]}
    \subbox{62.1mm}{\subh{2.2 典型时域性能指标}
    \[
      t_r=\frac{\pi-\theta}{\omega_d},\quad
      t_p=\frac{\pi}{\omega_d},\quad
      M_p=e^{-\pi\zeta/\sqrt{1-\zeta^2}}\times100\%
    \]
    \[
      t_s\approx\frac{4}{\zeta\omega_n}(2\%\text{准则}),\qquad
      t_s\approx\frac{3}{\zeta\omega_n}(5\%\text{准则})
    \]
    \centerline{\includegraphics[width=36mm,height=33mm,keepaspectratio]{step_response_crop.pdf}}
    }
    \subbox{62.1mm}{\subh{2.3 极点与稳定性}
    \begin{tabular}{|c|c|}
      \hline
      极点位置 & 系统性质\\ \hline
      全部左半平面 & 稳定\\
      含右半平面极点 & 不稳定\\
      虚轴单根 & 临界稳定\\ \hline
    \end{tabular}
    }
    \subbox{62.1mm}{\subh{2.4 终值与初值定理}
    \[
      f(0^+)=\lim_{s\to\infty}sF(s),\qquad
      f(\infty)=\lim_{s\to0}sF(s)
    \]}
    \subbox{62.1mm}{\subh{2.5 稳态误差}
    \[
      e_{ss}=\lim_{s\to0}sE(s),\quad
      E(s)=\frac{R(s)}{1+L(s)}
    \]
    \[
      K_p=\lim_{s\to0}L(s),\quad
      K_v=\lim_{s\to0}sL(s),\quad
      K_a=\lim_{s\to0}s^2L(s)
    \]
    \begin{tabular}{|c|c|c|c|}
      \hline
      型别 & 阶跃 & 斜坡 & 抛物线\\ \hline
      0型 & $1/(1+K_p)$ & $\infty$ & $\infty$\\
      1型 & $0$ & $1/K_v$ & $\infty$\\
      2型 & $0$ & $0$ & $1/K_a$\\ \hline
    \end{tabular}
    }
  }

  \contentnode{140.9}{35.2}{66.5mm}{
    \subbox{62.6mm}{\subh{3.1 定义}
    闭环特征方程 $1+KG(s)H(s)=0$。当 $K\in[0,\infty)$ 变化时，闭环极点在 $s$ 平面上的轨迹称为根轨迹。
    }
    \subbox{62.6mm}{\subh{3.2 根轨迹条件}
    \[
      \angle G(s)H(s)=(2k+1)\pi,\qquad
      K=\frac{1}{|G(s)H(s)|}
    \]}
    \subbox{62.6mm}{\subh{3.3 绘制规则}
    分支数等于开环极点数；关于实轴对称；实轴上右侧零极点总数为奇数的区段在轨迹上。
    \[
      \sigma_a=\frac{\sum p_i-\sum z_i}{n-m},\quad
      \theta_q=\frac{(2q+1)\pi}{n-m},\quad
      \frac{dK}{ds}=0
    \]}
    \subbox{62.6mm}{\subh{3.4 典型根轨迹示意}
    \centerline{\includegraphics[width=34mm,height=22mm,keepaspectratio]{root_locus_crop.pdf}}
    }
    \subbox{62.6mm}{\subh{3.5 主导极点}
    主导极点靠近虚轴时，对瞬态响应影响更强；轨迹左移通常提高快速性。
    }
  }

  \contentnode{140.9}{124.2}{66.5mm}{
    \subbox{62.6mm}{\subh{4.1 频率响应}
    \[
      G(j\omega)=|G(j\omega)|\angle G(j\omega),\qquad G(j\omega)=G(s)|_{s=j\omega}
    \]}
    \subbox{62.6mm}{\subh{4.2 Bode 基本要素}
    \begin{tabular}{|c|c|c|}
      \hline
      基本因子 & 幅值斜率 & 相位\\ \hline
      $K$ & $20\log K$ & $0^\circ$\\
      $1/s$ & $-20$dB/dec & $-90^\circ$\\
      $s$ & $+20$dB/dec & $+90^\circ$\\
      $Ts+1$ & $+20$dB/dec & $0^\circ\to90^\circ$\\
      $1/(Ts+1)$ & $-20$dB/dec & $0^\circ\to-90^\circ$\\ \hline
    \end{tabular}
    }
    \subbox{62.6mm}{\subh{4.3 Bode 图示意}
    \centerline{\includegraphics[width=37mm,height=32mm,keepaspectratio]{bode_plot_crop.pdf}}
    }
    \subbox{62.6mm}{\subh{4.4 闭环谐振}
    \[
      M_r=\frac{1}{2\zeta\sqrt{1-\zeta^2}},\quad
      \omega_r=\omega_n\sqrt{1-2\zeta^2}
    \]}
  }

  \contentnode{2.4}{183.2}{65mm}{
    \subbox{61.6mm}{\subh{5.1 常用校正网络}
    \begin{tabular}{|c|c|c|}
      \hline
      类型 & 传递函数 & 作用\\ \hline
      P & $K_p$ & 改变增益\\
      PI & $K_p(1+\frac1{T_is})$ & 降低稳态误差\\
      PD & $K_p(1+T_ds)$ & 改善阻尼\\
      PID & $K_p(1+\frac1{T_is}+T_ds)$ & 综合校正\\
      超前 & $K_c\frac{1+Ts}{1+\alpha Ts}$ & 增大相角裕度\\
      滞后 & $K_c\frac{1+Ts}{1+\beta Ts}$ & 提高低频增益\\ \hline
    \end{tabular}
    }
    \subbox{61.6mm}{\subh{5.2 超前/滞后零极点关系}
    \[
      \text{超前: }0<\alpha<1,\ |p|>|z|;\qquad
      \text{滞后: }\beta>1,\ |p|<|z|
    \]}
    \subbox{61.6mm}{\subh{5.3 设计流程}
    明确指标 $\rightarrow$ 选择校正类型 $\rightarrow$ 确定交越频率和相位补偿 $\rightarrow$ 校核 $e_{ss}$、$M_p$、$t_s$ 与稳定裕度。
    }
    \subbox{61.6mm}{\subh{5.4 调节经验}
    增大 $K$ 可加快响应但降低稳定裕度；积分改善稳态精度但可能增加振荡；微分改善阻尼但对噪声敏感。
    \[
      K\uparrow \Rightarrow \omega_c\uparrow,\quad
      \gamma\downarrow;\qquad
      \text{滞后校正常用于提高低频精度}
    \]}
  }

  \contentnode{71.4}{183.2}{66mm}{
    \subbox{62.1mm}{\subh{6.1 常用拉普拉斯变换}
    \begin{tabular}{|c|c||c|c|}
      \hline
      $f(t)$ & $F(s)$ & $f(t)$ & $F(s)$\\ \hline
      $1$ & $1/s$ & $t$ & $1/s^2$\\
      $t^2/2$ & $1/s^3$ & $e^{-at}$ & $1/(s+a)$\\
      $\sin\omega t$ & $\omega/(s^2+\omega^2)$ & $\cos\omega t$ & $s/(s^2+\omega^2)$\\ \hline
    \end{tabular}
    }
    \subbox{62.1mm}{\subh{6.2 经典结论}
    1) 型别由 $L(s)$ 在原点的极点数决定。\\
    2) 低频增益主导稳态误差，高频衰减影响噪声抑制。\\
    3) 主导共轭极点决定欠阻尼响应的速度与超调。\\
    4) 非最小相位零点引入额外相位滞后，整定难度增加。\\
    5) 根轨迹向左通常提高快速性，但仍需检查阻尼比。\\
    6) 相位裕度过小会导致振荡敏感，过大则响应可能偏慢。
    }
    \subbox{62.1mm}{\subh{6.3 部分分式}
    \[
      \frac{N(s)}{D(s)}
      =\sum_i\frac{A_i}{(s-p_i)^r}
      +\sum_j\frac{B_js+C_j}{s^2+2\zeta_j\omega_js+\omega_j^2}
    \]
    常用判断：闭环极点决定自然响应，开环型别决定典型输入下的稳态误差，频率裕度衡量对建模误差的容忍度。
    }
  }

  \contentnode{140.9}{202.2}{66.5mm}{
    \subbox{62.6mm}{\subh{4.6 奈奎斯特判据}
    \[
      N=Z-P,\qquad Z=0\Rightarrow N=-P
    \]
    \centerline{\includegraphics[width=28mm,height=21mm,keepaspectratio]{nyquist_plot_crop.pdf}}
    }
    \subbox{62.6mm}{\subh{4.7 稳定裕度}
    \[
      \omega_c:\ |L(j\omega_c)|=1,\quad
      \omega_\pi:\angle L(j\omega_\pi)=-180^\circ
    \]
    \[
      \gamma=180^\circ+\angle L(j\omega_c),\quad
      G_m=-20\log_{10}|L(j\omega_\pi)|
    \]}
    \subbox{62.6mm}{\subh{4.8 经验判断}
    $\gamma>60^\circ$ 阻尼较强；$30^\circ\sim60^\circ$ 常用；$0^\circ\sim30^\circ$ 振荡敏感。
    }
  }

  \node[anchor=north west,inner sep=.8mm,text width=50mm] at (2.4,270)
    {\boxfont\subh{7.1 阶跃响应}\includegraphics[width=30mm,height=17mm,keepaspectratio]{step_response_crop.pdf}};
  \node[anchor=north west,inner sep=.8mm,text width=50mm] at (54,270)
    {\boxfont\subh{7.2 根轨迹图}\includegraphics[width=30mm,height=17mm,keepaspectratio]{root_locus_crop.pdf}};
  \node[anchor=north west,inner sep=.8mm,text width=50mm] at (105.5,270)
    {\boxfont\subh{7.3 Bode 图}\includegraphics[width=30mm,height=17mm,keepaspectratio]{bode_plot_crop.pdf}};
  \node[anchor=north west,inner sep=.8mm,text width=50mm] at (157,270)
    {\boxfont\subh{7.4 Nyquist 图}\includegraphics[width=30mm,height=17mm,keepaspectratio]{nyquist_plot_crop.pdf}};
  \draw[ruleblue,line width=.25pt] (53.2,269.5) -- (53.2,291);
  \draw[ruleblue,line width=.25pt] (105,269.5) -- (105,291);
  \draw[ruleblue,line width=.25pt] (156.5,269.5) -- (156.5,291);

  \node[anchor=west,font=\fontsize{4.8pt}{5.2pt}\selectfont] at (2,294.2)
    {注：本页为经典控制理论矢量电子版速查表，公式、表格、框图与曲线图可放大阅读；工程使用时需结合系统结构与指标约束校核。};
\end{tikzpicture}
\end{document}
